\section{Preliminaries}
\label{sec:preliminaries}
% \subsection{Notation}
% \label{sec:notation}
\textbf{Notation.} We use $\langle\cdot,\cdot\rangle$ and $\|\cdot\|$ for the
Euclidean inner product and norm; for a matrix, $\|\cdot\|$ is the spectral
norm. For $\calA=(A_1,\ldots,A_m)^\top$, the Jacobian
$J_{\calA}(\bx)\in\RR^{m\times n}$ has $i$th row $\nabla A_i(\bx)^\top$. For a
nonempty closed convex set $C$, $\dist(\bu,C):=\min_{\bv\in C}\|\bu-\bv\|$, and
$\Pi_C(\bu)$ is the Euclidean projection of $\bu$ onto $C$. When convenient, we
identify $\RR^2$ with $\mathbb C$, so that $\|p\|=|p|$ and
$\langle p,q\rangle=\operatorname{Re}(\bar pq)$.
% Throughout the paper, $\partial h$ denotes the limiting
% subdifferential of $h$.

For a proper closed convex function \(h:\RR^m\to\RR\cup\{+\infty\}\), its
subdifferential, convex conjugate, and proximal mapping are
\[
    \partial h(\by):=\{\bv:\ h(\bw)\ge h(\by)+\langle\bv,\bw-\by\rangle
    \ \text{for all }\bw\},
    \qquad
    h^*(\bq):=\sup_{\by\in\RR^m}\{\langle\bq,\by\rangle-h(\by)\},
\]
\begin{equation}
    \prox_{t h}(\bu)
    :=\argmin_{\by\in\RR^m}
       \Bigl\{h(\by)+\frac{1}{2t}\|\by-\bu\|^2\Bigr\},
    \qquad t>0;
    \label{eq:prox-definition}
\end{equation}
see~\cite{rockafellar1970convex}. We write
\(\dom\partial h^*:=\{\bq:\partial h^*(\bq)\neq\varnothing\}\).


\subsection{Assumptions}
\label{sec:assumptions}
We impose the following assumptions.
% \begin{assumption}[Uniform Jacobian regularity]
% \label{ass:strong-licq}
% There exists a constant $\nu>0$ such that
% \[
%     \sigma_{\min}\bigl(J_{\calA}(\bx)\bigr)\ge \nu,
%     \qquad
%     \forall \bx\in\calS.
% \]
% Equivalently,
% \[
%     \|J_{\calA}(\bx)^\top \by\|
%     \ge
%     \nu \|\by\|,
%     \qquad \forall \bx\in\calS,\ \forall \by\in\RR^m.
% \]
% \end{assumption}
\begin{assumption}[Smooth components]
\label{ass:smoothness}
There exist constants $L_f,G_f,L_{\calA},G_{\calA}\ge0$ such that, for
all $\bx,\bx'\in\RR^n$,
\[
\begin{aligned}
    \|\nabla f(\bx)-\nabla f(\bx')\|
    &\le L_f\|\bx-\bx'\|,
    &\|\nabla f(\bx)\|&\le G_f,\\
    \|J_{\calA}(\bx)-J_{\calA}(\bx')\|
    &\le L_{\calA}\|\bx-\bx'\|,
    &\|J_{\calA}(\bx)\|&\le G_{\calA}.
\end{aligned}
\]
\end{assumption}

\begin{assumption}[Lipschitz outer function]
\label{ass:h}
The function $h:\RR^m\to\RR$ is convex and $L_h$-Lipschitz:
\[
    |h(\bu)-h(\bv)|\le L_h\|\bu-\bv\|,
    \qquad \bu,\bv\in\RR^m.
\]
\end{assumption}

\subsection{Problem class, stationarity, and oracle}
\label{subsec:setup}
To study~(Q), we introduce the following problem class.
\begin{definition}[Problem class]
\label{def:problem-class}
Let \(\Delta>0\) and
\(\mathbf C=(L_f,G_f,L_{\calA},G_{\calA},L_h)\in\RR^5_{\ge0}\). 
We denote by \(\mathcal I(\mathbf C,\Delta)\) the set of all triples
\((f,\calA,h)\), over all dimensions \(n,m\in\NN\), such that
Assumptions~\ref{ass:smoothness} and~\ref{ass:h} hold with the constants
in \(\mathbf C\) and 
\(
    F(\zero)-\inf_{\bx\in\RR^n}F(\bx)\le\Delta .
\)
\end{definition}
The constants \(\mathbf C\) and \(\Delta\) do not depend on the target
accuracy or on \(n,m\), and the gap condition implies \(\inf F>-\infty\).
% The gap condition implies 
% \(\inf F>-\infty\). Assumption~\ref{ass:level-boundedness} is not part of
% \(\mathcal I(\mathbf C,\Delta)\); it is imposed separately in the complexity
% results that require it. 

{We use the following stationarity measure for Problem~\eqref{prob:main}.} 
\begin{definition}[\(\epsilon\)-KKT point]
\label{def:eps-kkt}
Given \(\epsilon>0\). A triple \((\bx,\by,\bq)\) with
\(\bq\in\partial h(\by)\) is an \(\epsilon\)-KKT point of
\eqref{prob:main} if
\begin{equation}
    \mathcal{R}_{\mathrm{KKT}}(\bx,\by,\bq)
    :=
    \max\Bigl\{
        \bigl\|\nabla f(\bx)+J_{\calA}(\bx)^\top\bq\bigr\|,\;
        \|\calA(\bx)-\by\|
    \Bigr\} \leq \epsilon
    \label{eq:kkt-residual}
\end{equation}
\end{definition}

We next relate Definition~\ref{def:eps-kkt} to game stationarity (GS) for the
minimax reformulation~Problem~\eqref{eq:ncl} of Problem~\eqref{prob:main}. We use the
following GS definition from \cite[Definition~3.1(b)]{li2025nonsmooth} and
\cite[(2.10)--(2.11)]{zhang2026matching}.
\begin{definition}[\(\epsilon\)-GS point]
\label{def:eps-gs}
Let \(\epsilon>0\). A pair \((\bx,\bq)\in\RR^n\times\dom\partial h^*\) is an
\(\epsilon\)-game stationary point (\(\epsilon\)-GS) of~Problem~\eqref{eq:ncl} if
\[
    \calG(\bx,\bq)
    :=
    \max\Bigl\{
        \bigl\|\nabla f(\bx)+J_{\calA}(\bx)^\top\bq\bigr\|,\;
        \dist\bigl(\calA(\bx),\partial h^*(\bq)\bigr)
    \Bigr\}
    \le\epsilon .
\]
\end{definition}
We now relate $\epsilon$-KKT points of Problem~\eqref{prob:main}
to $\epsilon$-GS points of its Fenchel saddle reformulation~Problem~\eqref{eq:ncl}. 
\begin{proposition}[KKT points of Problem~\eqref{prob:main} and GS points of~Problem~\eqref{eq:ncl}]
\label{prop:kkt-gs}
Suppose that Assumption~\ref{ass:h} holds, and let
\((\bx,\bq)\in\RR^n\times\dom\partial h^*\). Then the following measurement relationship holds: 
\[
    \calG(\bx,\bq)
    =
    \min_{\by\in\partial h^*(\bq)}
    \mathcal{R}_{\mathrm{KKT}}(\bx,\by,\bq),
\]
with minimum attained at
\(\hat\by=\Pi_{\partial h^*(\bq)}\bigl(\calA(\bx)\bigr)\). Consequently, for
every \(\epsilon>0\):
\begin{enumerate}[label=(\alph*)]
\item if \((\bx,\by,\bq)\) is an \(\epsilon\)-KKT point
  of~Problem~\eqref{prob:main}, then \((\bx,\bq)\) is an \(\epsilon\)-GS point
  of~Problem~\eqref{eq:ncl};
\item if \((\bx,\bq)\) is an \(\epsilon\)-GS point of~Problem~\eqref{eq:ncl}, then
  \((\bx,\hat\by,\bq)\) is an \(\epsilon\)-KKT point of~Problem~\eqref{prob:main}.
\end{enumerate}
\end{proposition}

\begin{proof}
Since \(h\) is proper, closed, and convex,
\(\bq\in\partial h(\by)\) if and only if \(\by\in\partial h^*(\bq)\)
\cite[Theorem~23.5]{rockafellar1970convex}. The first term of
\(\mathcal{R}_{\mathrm{KKT}}\) does not depend on \(\by\). The set
\(\partial h^*(\bq)\) is nonempty, closed, and convex, so
\(\min_{\by\in\partial h^*(\bq)}\|\calA(\bx)-\by\|
=\dist(\calA(\bx),\partial h^*(\bq))\), attained at \(\hat\by\).
Statements~(a) and~(b) follow.
\end{proof}

We next introduce the following deterministic first-order oracle model.
\begin{definition}[Oracle model]
\label{def:oracle-model}
We assume access to a deterministic first-order oracle
\[
    \mathcal{O}_{f,\mathcal{A}}(\bx)
    :=
    \bigl(f(\bx),\nabla f(\bx),\mathcal{A}(\bx),J_{\mathcal{A}}(\bx)\bigr),
\]
and an exact proximal oracle
\[
    \mathcal{O}_{h}(\bv,t)
    :=
    \operatorname{prox}_{t h}(\bv),
    \qquad t>0.
\]
Each call to either oracle counts as one oracle evaluation. In particular,
the cost of solving the proximal subproblem is represented by one exact
proximal-oracle call.
\end{definition}

